% =====================================================
% 题型：等腰直角三角形中平面向量综合多选题 (q_vector_triangle_intersection.tex)
% =====================================================
% 难度：★★ (中档)
% =====================================================
\newcommand{\qVectorRightTriangleIntersectionStem}{在 \(\triangle ABC\) 中，已知 \(\angle BAC=\dfrac{\pi}{2}\)，\(AB=AC\)，\(D\) 是 \(AC\) 的中点。若 \(P\) 是 \(BC\) 上的一点，且满足 \(\overrightarrow{BP}=2\overrightarrow{PC}\)，\(AP\) 与 \(BD\) 交于点 \(E\)，则下列结论正确的是 \hfill （\quad）}

\newcommand{\qVectorRightTriangleIntersectionOptions}{%
  \mcoptions[1]{\(\overrightarrow{AP} = \dfrac13\overrightarrow{AB} + \dfrac23\overrightarrow{AC}\)}%
               {\(\overrightarrow{AP}\) 在 \(\overrightarrow{AB}\) 上的投影向量为 \(\dfrac23\overrightarrow{AB}\)}%
               {\(\overrightarrow{AP}\cdot\overrightarrow{BD}=0\)}%
               {\(\overrightarrow{AE}=\dfrac35\overrightarrow{AP}\)}%
}

\newcommand{\qVectorRightTriangleIntersectionAnswer}{ACD}

\newcommand{\qVectorRightTriangleIntersectionFigure}[1][0.85]{%
  \begin{tikzpicture}[scale=#1, line join=round, line cap=round, >=stealth]
    % Coordinates: A(0,0), B(3,0), C(0,3)
    \coordinate (A) at (0,0);
    \coordinate (B) at (3,0);
    \coordinate (C) at (0,3);
    \coordinate (D) at (0,1.5);
    \coordinate (P) at (1,2);
    \coordinate (E) at (0.6,1.2);

    % Draw triangle
    \draw[very thick] (A) -- (B) -- (C) -- cycle;

    % Draw lines AP and BD
    \draw[thick] (A) -- (P);
    \draw[thick] (B) -- (D);

    % Nodes
    \fill (A) circle (1.2pt) node[below left] {\small \(A\)};
    \fill (B) circle (1.2pt) node[below right] {\small \(B\)};
    \fill (C) circle (1.2pt) node[above left] {\small \(C\)};
    \fill (D) circle (1.2pt) node[left] {\small \(D\)};
    \fill (P) circle (1.2pt) node[above right] {\small \(P\)};
    \fill (E) circle (1.2pt) node[below left, xshift=-1pt, yshift=-1pt] {\small \(E\)};

    % Right angle sign at A
    \draw (0.2,0) -- (0.2,0.2) -- (0,0.2);

    \ifteacher
      % Perpendicular sign at E
      % AE direction angle is approx 63.4 degrees
      \begin{scope}[shift={(E)}, rotate=63.4]
        \draw[gray, thin] (0.18,0) -- (0.18,0.18) -- (0,0.18);
      \end{scope}
    \fi
  \end{tikzpicture}%
}

\newcommand{\qVectorRightTriangleIntersectionSolution}{%
  设 \(\overrightarrow{AB}=\vec{u}\)，\(\overrightarrow{AC}=\vec{v}\)。
  因为 \(\angle BAC=\dfrac{\pi}{2}\)，且 \(AB=AC\)，所以 \(\vec{u} \cdot \vec{v} = 0\)，且 \(|\vec{u}| = |\vec{v}|\)。

  ---

  【对于 A 项】
  因为 \(\overrightarrow{BP}=2\overrightarrow{PC}\)，根据分点公式可得：
  \[
    \overrightarrow{BP} = \frac{2}{3}\overrightarrow{BC}.
  \]
  而 \(\overrightarrow{BC} = \overrightarrow{AC} - \overrightarrow{AB} = \vec{v} - \vec{u}\)，
  所以：
  \[
    \overrightarrow{AP} = \overrightarrow{AB} + \overrightarrow{BP} = \vec{u} + \frac{2}{3}(\vec{v} - \vec{u}) = \frac{1}{3}\vec{u} + \frac{2}{3}\vec{v} = \frac{1}{3}\overrightarrow{AB} + \frac{2}{3}\overrightarrow{AC}.
  \]
  故 A 正确。

  ---

  【对于 B 项】
  \(\overrightarrow{AP}\) 在 \(\overrightarrow{AB}\) 上的投影向量为：
  \[
    \operatorname{proj}_{\vec{u}}\overrightarrow{AP} = \frac{\overrightarrow{AP} \cdot \vec{u}}{|\vec{u}|^2} \vec{u} = \frac{\left(\frac{1}{3}\vec{u} + \frac{2}{3}\vec{v}\right) \cdot \vec{u}}{|\vec{u}|^2} \vec{u}.
  \]
  因为 \(\vec{u} \cdot \vec{v} = 0\)，所以上式为：
  \[
    \frac{\frac{1}{3}|\vec{u}|^2}{|\vec{u}|^2} \vec{u} = \frac{1}{3}\vec{u} = \frac{1}{3}\overrightarrow{AB}.
  \]
  故 B 错误。

  ---

  【对于 C 项】
  因为 \(D\) 是 \(AC\) 的中点，所以 \(\overrightarrow{AD} = \dfrac{1}{2}\overrightarrow{AC} = \dfrac{1}{2}\vec{v}\)。
  进而：
  \[
    \overrightarrow{BD} = \overrightarrow{AD} - \overrightarrow{AB} = \frac{1}{2}\vec{v} - \vec{u}.
  \]
  计算数量积 \(\overrightarrow{AP} \cdot \overrightarrow{BD}\)：
  \[
    \overrightarrow{AP} \cdot \overrightarrow{BD} = \left(\frac{1}{3}\vec{u} + \frac{2}{3}\vec{v}\right) \cdot \left(\frac{1}{2}\vec{v} - \vec{u}\right) = -\frac{1}{3}\vec{u}^2 + \frac{1}{3}\vec{v}^2 + \left(\frac{1}{6} - \frac{2}{3}\right)\vec{u} \cdot \vec{v}.
  \]
  由于 \(\vec{u} \cdot \vec{v} = 0\) 且 \(|\vec{u}| = |\vec{v}|\)，因此：
  \[
    \overrightarrow{AP} \cdot \overrightarrow{BD} = -\frac{1}{3}|\vec{u}|^2 + \frac{1}{3}|\vec{v}|^2 = 0 \implies AP \perp BD.
  \]
  故 C 正确。

  ---

  【对于 D 项】
  因为 \(E\) 是直线 \(AP\) 与 \(BD\) 的交点。
  设 \(\overrightarrow{AE} = \lambda \overrightarrow{AP} = \dfrac{\lambda}{3}\vec{u} + \dfrac{2\lambda}{3}\vec{v}\)。

  又因为 \(E\) 在直线 \(BD\) 上，根据共线向量定理，存在实数 \(\mu\) 使得：
  \[
    \overrightarrow{AE} = \overrightarrow{AB} + \mu \overrightarrow{BD} = \vec{u} + \mu\left(\frac{1}{2}\vec{v} - \vec{u}\right) = (1 - \mu)\vec{u} + \frac{\mu}{2}\vec{v}.
  \]

  由于基底 \(\vec{u}, \vec{v}\) 分解唯一，比较系数得：
  \[
    \begin{cases}
      \dfrac{\lambda}{3} = 1 - \mu \\[6pt]
      \dfrac{2\lambda}{3} = \dfrac{\mu}{2}
    \end{cases}
    \implies \mu = \frac{4\lambda}{3}.
  \]
  代入第一式：
  \[
    \frac{\lambda}{3} = 1 - \frac{4\lambda}{3} \implies \frac{5\lambda}{3} = 1 \implies \lambda = \frac{3}{5}.
  \]
  所以 \(\overrightarrow{AE} = \dfrac{3}{5}\overrightarrow{AP}\)。
  故 D 正确。

  ---

  综上所述，结论正确的是 A, C, D。
  故选 ACD。%
}
